% Options for packages loaded elsewhere
\PassOptionsToPackage{unicode}{hyperref}
\PassOptionsToPackage{hyphens}{url}
%
\documentclass[
]{article}
\usepackage{amsmath,amssymb}
\usepackage{iftex}
\ifPDFTeX
  \usepackage[T1]{fontenc}
  \usepackage[utf8]{inputenc}
  \usepackage{textcomp} % provide euro and other symbols
\else % if luatex or xetex
  \usepackage{unicode-math} % this also loads fontspec
  \defaultfontfeatures{Scale=MatchLowercase}
  \defaultfontfeatures[\rmfamily]{Ligatures=TeX,Scale=1}
\fi
\usepackage{lmodern}
\ifPDFTeX\else
  % xetex/luatex font selection
\fi
% Use upquote if available, for straight quotes in verbatim environments
\IfFileExists{upquote.sty}{\usepackage{upquote}}{}
\IfFileExists{microtype.sty}{% use microtype if available
  \usepackage[]{microtype}
  \UseMicrotypeSet[protrusion]{basicmath} % disable protrusion for tt fonts
}{}
\makeatletter
\@ifundefined{KOMAClassName}{% if non-KOMA class
  \IfFileExists{parskip.sty}{%
    \usepackage{parskip}
  }{% else
    \setlength{\parindent}{0pt}
    \setlength{\parskip}{6pt plus 2pt minus 1pt}}
}{% if KOMA class
  \KOMAoptions{parskip=half}}
\makeatother
\usepackage{xcolor}
\usepackage[margin=1in]{geometry}
\usepackage{color}
\usepackage{fancyvrb}
\newcommand{\VerbBar}{|}
\newcommand{\VERB}{\Verb[commandchars=\\\{\}]}
\DefineVerbatimEnvironment{Highlighting}{Verbatim}{commandchars=\\\{\}}
% Add ',fontsize=\small' for more characters per line
\usepackage{framed}
\definecolor{shadecolor}{RGB}{248,248,248}
\newenvironment{Shaded}{\begin{snugshade}}{\end{snugshade}}
\newcommand{\AlertTok}[1]{\textcolor[rgb]{0.94,0.16,0.16}{#1}}
\newcommand{\AnnotationTok}[1]{\textcolor[rgb]{0.56,0.35,0.01}{\textbf{\textit{#1}}}}
\newcommand{\AttributeTok}[1]{\textcolor[rgb]{0.13,0.29,0.53}{#1}}
\newcommand{\BaseNTok}[1]{\textcolor[rgb]{0.00,0.00,0.81}{#1}}
\newcommand{\BuiltInTok}[1]{#1}
\newcommand{\CharTok}[1]{\textcolor[rgb]{0.31,0.60,0.02}{#1}}
\newcommand{\CommentTok}[1]{\textcolor[rgb]{0.56,0.35,0.01}{\textit{#1}}}
\newcommand{\CommentVarTok}[1]{\textcolor[rgb]{0.56,0.35,0.01}{\textbf{\textit{#1}}}}
\newcommand{\ConstantTok}[1]{\textcolor[rgb]{0.56,0.35,0.01}{#1}}
\newcommand{\ControlFlowTok}[1]{\textcolor[rgb]{0.13,0.29,0.53}{\textbf{#1}}}
\newcommand{\DataTypeTok}[1]{\textcolor[rgb]{0.13,0.29,0.53}{#1}}
\newcommand{\DecValTok}[1]{\textcolor[rgb]{0.00,0.00,0.81}{#1}}
\newcommand{\DocumentationTok}[1]{\textcolor[rgb]{0.56,0.35,0.01}{\textbf{\textit{#1}}}}
\newcommand{\ErrorTok}[1]{\textcolor[rgb]{0.64,0.00,0.00}{\textbf{#1}}}
\newcommand{\ExtensionTok}[1]{#1}
\newcommand{\FloatTok}[1]{\textcolor[rgb]{0.00,0.00,0.81}{#1}}
\newcommand{\FunctionTok}[1]{\textcolor[rgb]{0.13,0.29,0.53}{\textbf{#1}}}
\newcommand{\ImportTok}[1]{#1}
\newcommand{\InformationTok}[1]{\textcolor[rgb]{0.56,0.35,0.01}{\textbf{\textit{#1}}}}
\newcommand{\KeywordTok}[1]{\textcolor[rgb]{0.13,0.29,0.53}{\textbf{#1}}}
\newcommand{\NormalTok}[1]{#1}
\newcommand{\OperatorTok}[1]{\textcolor[rgb]{0.81,0.36,0.00}{\textbf{#1}}}
\newcommand{\OtherTok}[1]{\textcolor[rgb]{0.56,0.35,0.01}{#1}}
\newcommand{\PreprocessorTok}[1]{\textcolor[rgb]{0.56,0.35,0.01}{\textit{#1}}}
\newcommand{\RegionMarkerTok}[1]{#1}
\newcommand{\SpecialCharTok}[1]{\textcolor[rgb]{0.81,0.36,0.00}{\textbf{#1}}}
\newcommand{\SpecialStringTok}[1]{\textcolor[rgb]{0.31,0.60,0.02}{#1}}
\newcommand{\StringTok}[1]{\textcolor[rgb]{0.31,0.60,0.02}{#1}}
\newcommand{\VariableTok}[1]{\textcolor[rgb]{0.00,0.00,0.00}{#1}}
\newcommand{\VerbatimStringTok}[1]{\textcolor[rgb]{0.31,0.60,0.02}{#1}}
\newcommand{\WarningTok}[1]{\textcolor[rgb]{0.56,0.35,0.01}{\textbf{\textit{#1}}}}
\usepackage{graphicx}
\makeatletter
\def\maxwidth{\ifdim\Gin@nat@width>\linewidth\linewidth\else\Gin@nat@width\fi}
\def\maxheight{\ifdim\Gin@nat@height>\textheight\textheight\else\Gin@nat@height\fi}
\makeatother
% Scale images if necessary, so that they will not overflow the page
% margins by default, and it is still possible to overwrite the defaults
% using explicit options in \includegraphics[width, height, ...]{}
\setkeys{Gin}{width=\maxwidth,height=\maxheight,keepaspectratio}
% Set default figure placement to htbp
\makeatletter
\def\fps@figure{htbp}
\makeatother
\setlength{\emergencystretch}{3em} % prevent overfull lines
\providecommand{\tightlist}{%
  \setlength{\itemsep}{0pt}\setlength{\parskip}{0pt}}
\setcounter{secnumdepth}{5}
\usepackage{booktabs}
\usepackage{longtable}
\usepackage{array}
\usepackage{multirow}
\usepackage{wrapfig}
\usepackage{float}
\usepackage{colortbl}
\usepackage{pdflscape}
\usepackage{tabu}
\usepackage{threeparttable}
\usepackage{threeparttablex}
\usepackage[normalem]{ulem}
\usepackage{makecell}
\usepackage{xcolor}
\ifLuaTeX
  \usepackage{selnolig}  % disable illegal ligatures
\fi
\usepackage{bookmark}
\IfFileExists{xurl.sty}{\usepackage{xurl}}{} % add URL line breaks if available
\urlstyle{same}
\hypersetup{
  pdftitle={Analyse des établissements cinématographiques en France},
  pdfauthor={Axel Frache, Nathan Dilhan, Liam Soulet},
  hidelinks,
  pdfcreator={LaTeX via pandoc}}

\title{Analyse des établissements cinématographiques en France}
\author{Axel Frache, Nathan Dilhan, Liam Soulet}
\date{2025-04-01}

\begin{document}
\maketitle

{
\setcounter{tocdepth}{2}
\tableofcontents
}
\section{Introduction}\label{introduction}

Ce rapport propose une analyse statistique descriptive des
établissements cinématographiques en France à partir d'un jeu de données
riche de 2061 observations et 40 variables.

Nous cherchons à répondre aux questions suivantes :

\begin{itemize}
\tightlist
\item
  \textbf{Comment les cinémas sont-ils répartis sur le territoire
  français ?}
\item
  \textbf{Existe-t-il des différences structurelles entre les
  établissements (nombre d'écrans, capacité, type) ?}
\item
  \textbf{Quel est le lien entre les caractéristiques d'un établissement
  (type, capacité) et sa fréquentation ?}
\item
  \textbf{Comment se répartit la part de marché du cinéma français et du
  cinéma Art et Essai ?}
\end{itemize}

À travers des visualisations claires, nous dresserons un état des lieux
synthétique du paysage cinématographique français.

\section{Description du jeu de
données}\label{description-du-jeu-de-donnuxe9es}

Le jeu de données contient des informations sur les cinémas français :
localisation, équipements (écrans, fauteuils), fréquentation (entrées,
séances), typologie (multiplexe, art et essai), et part de marché par
type de films. Voici un aperçu des premières lignes :

\begin{longtable}[t]{rrlllllrlllrlrrrrrrrllllllllrrrrrrrrrrrl}
\toprule
régionCNC & N° auto & nom & région administrative & adresse & code INSEE & commune & population de la commune & DEP & N°UU & unité urbaine & population unité urbaine & situation géographique & écrans & fauteuils & semaines d'activité & séances & entrées 2022 & entrées 2021 & évolution entrées & tranche d'entrées & programmateur & AE & catégorie Art et Essai & label Art et Essai & genre & multiplexe & zone de la commune & nombre de films programmés & nombre de films inédits & nombre de films en semaine 1 & PdM en entrées des films français & PdM en entrées des films américains & PdM en entrées des films européens & PdM en entrées des autres films & films Art et Essai & PdM en entrées des films Art et Essai & latitude & longitude & geolocalisation\\
\midrule
1 & 31 & UGC NORMANDIE & ILE-DE-FRANCE & 116 AVENUE DES CHAMPS ELYSEES & 75108 & Paris 8e Arrondissement & 36218 & 75 & 00851 & Paris & 10949158 & 01-Paris & 4 & 1533 & 53 & 6966 & 148007 & 119434 & 23.92367 & De 100 à 150 000 entrées & UGC & NON & NON & NA & FIXE & NON & C & 125 & 79 & 77 & 44.70344 & 36.505312 & 12.45128 & 6.339970 & 54 & 23.97523 & 48.87226 & 2.300938 & 48.872265, 2.300938\\
1 & 204 & UGC OPERA & ILE-DE-FRANCE & 32 BD DES ITALIENS & 75109 & Paris 9e Arrondissement & 60784 & 75 & 00851 & Paris & 10949158 & 01-Paris & 4 & 722 & 53 & 6840 & 110063 & 81837 & 34.49051 & De 100 à 150 000 entrées & UGC & NON & NON & NA & FIXE & NON & C & 111 & 71 & 71 & 55.64449 & 27.053597 & 12.62277 & 4.679138 & 55 & 23.24759 & 48.87133 & 2.335302 & 48.871332, 2.335302\\
1 & 733 & STUDIO GALANDE & ILE-DE-FRANCE & 42 RUE GALANDE & 75105 & Paris 5e Arrondissement & 58050 & 75 & 00851 & Paris & 10949158 & 01-Paris & 1 & 97 & 53 & 996 & 13643 & 8084 & 68.76546 & De 10 à 20 000 entrées & NA & OUI & A & NA & FIXE & NON & C & 155 & 110 & 4 & 25.68362 & 4.697148 & 53.33725 & 16.281976 & 115 & 87.83446 & 48.85162 & 2.347076 & 48.851619, 2.347076\\
1 & 821 & PARNASSIEN & ILE-DE-FRANCE & 98 BD DU MONTPARNASSE & 75114 & Paris 14e Arrondissement & 134926 & 75 & 00851 & Paris & 10949158 & 01-Paris & 7 & 842 & 53 & 13507 & 216828 & 95714 & 126.53739 & De 200 à 300 000 entrées & NA & OUI & A & RD & FIXE & NON & C & 248 & 153 & 134 & 63.36218 & 5.584636 & 15.32115 & 15.732034 & 201 & 94.58166 & 48.84244 & 2.327625 & 48.842443, 2.327625\\
1 & 1031 & PUBLICIS ELYSEES & ILE-DE-FRANCE & 133 AV DES CHAMPS ELYSEES & 75108 & Paris 8e Arrondissement & 36218 & 75 & 00851 & Paris & 10949158 & 01-Paris & 2 & 609 & 50 & 3164 & 46020 & 35847 & 28.37894 & De 40 à 50 000 entrées & NA & NON & NON & NA & FIXE & NON & C & 165 & 86 & 48 & 38.38598 & 20.324249 & 16.80629 & 24.483483 & 68 & 42.59715 & 48.87279 & 2.296984 & 48.872794, 2.296984\\
\bottomrule
\end{longtable}

\section{Analyse descriptive}\label{analyse-descriptive}

\subsection{Répartition des établissements par
région}\label{ruxe9partition-des-uxe9tablissements-par-ruxe9gion}

Cette première analyse permet de répondre à la question : \textbf{Quels
territoires sont les mieux équipés en salles de cinéma ?}

\includegraphics{rapport_files/figure-latex/region-repartition-1.pdf}

\begin{quote}
\textbf{Observation :} L'Île-de-France concentre naturellement de
nombreux établissements, mais certaines régions comme
l'Auvergne-Rhône-Alpes ou PACA sont aussi bien représentées. Cela laisse
penser que l'implantation des cinémas ne dépend pas uniquement de la
population qui y réside, mais aussi de la fréquentation temporaire
(tourisme, festivals\ldots).
\end{quote}

\subsection{Distribution du nombre d'écrans par
établissement}\label{distribution-du-nombre-duxe9crans-par-uxe9tablissement}

Cette analyse s'intéresse à la structure des établissements :
\textbf{sont-ils majoritairement de petite taille ou de grands complexes
?}

\includegraphics{rapport_files/figure-latex/histogramme-ecrans-1.pdf}

\begin{quote}
\textbf{Observation :} Ce qu'on remarque surtout ici, c'est que la
plupart des cinémas en France sont de petite taille. C'est souvent des
cinés de centre-ville, parfois indépendants ou associatifs, avec juste
quelques écrans. Les gros multiplexes (avec 8 ou 10 salles) existent,
mais ils sont clairement pas la majorité. C'est plutôt cohérent avec le
fait que la France soutient pas mal ses petites salles, notamment via
les aides au cinéma Art et Essai.
\end{quote}

\subsection{Entrées selon le type d'établissement
(multiplexe)}\label{entruxe9es-selon-le-type-duxe9tablissement-multiplexe}

Nous étudions ici : \textbf{le type d'établissement influence-t-il la
fréquentation ?}

\includegraphics{rapport_files/figure-latex/boxplot-multiplexe-1.pdf}

\begin{quote}
\textbf{Observation :} Le boxplot montre clairement un écart de
fréquentation selon le type d'établissement. Les multiplexes concentrent
l'essentiel des entrées, ce qui s'explique par leur capacité d'accueil
plus grande et une programmation souvent plus commerciale. Cependant, le
fait que certains petits cinémas ``classiques'' atteignent aussi des
niveaux élevés de fréquentation (en particulier sur les extrêmes du
boxplot) suggère qu'il existe des contre-exemples : peut-être des
établissements situés dans des zones sans concurrence ou qui bénéficient
d'une vraie fidélité du public local.
\end{quote}

\subsection{Corrélation entre le nombre de fauteuils et les
entrées}\label{corruxe9lation-entre-le-nombre-de-fauteuils-et-les-entruxe9es}

Nous explorons ici la question : \textbf{la capacité d'accueil est-elle
un facteur explicatif de la fréquentation ?}

\includegraphics{rapport_files/figure-latex/correlation-fauteuils-1.pdf}

\begin{quote}
\textbf{Observation :} Une relation linéaire est visible : plus un
établissement dispose de places, plus le volume de spectateurs est
élevé, bien que la corrélation ne soit pas parfaite. On observe quand
même une certaine dispersion : à nombre de fauteuils égal, la
fréquentation peut fortement varier. Cela montre que d'autres facteurs
entrent en jeu (type de programmation, zone géographique, réputation du
cinéma\ldots). En gros, avoir plus de places aide, mais ça ne fait pas
tout.
\end{quote}

\subsection{Part de marché des films
français}\label{part-de-marchuxe9-des-films-franuxe7ais}

Nous cherchons à visualiser : \textbf{quelle place occupe le cinéma
français dans la programmation ?}

\includegraphics{rapport_files/figure-latex/pdm-francais-1.pdf}

\begin{quote}
\textbf{Observation :} Le graphe montre une présence forte du cinéma
français, ce qui n'est pas anodin. Dans pas mal d'établissements, la
part des entrées générées par les films français dépasse les 40 \%, ce
qui témoigne de leur bonne diffusion. C'est probablement lié à plusieurs
facteurs : politique de quotas, attachement culturel, et aussi une
production assez diversifiée qui permet de toucher des publics variés.
\end{quote}

\subsection{Art et Essai selon zone
géographique}\label{art-et-essai-selon-zone-guxe9ographique}

Dernière question abordée : \textbf{la localisation géographique
influence-t-elle la programmation culturelle (art et essai) ?}

\includegraphics{rapport_files/figure-latex/art-essai-zone-1.pdf}

\begin{quote}
\textbf{Observation :} Ce résultat est assez révélateur : les cinémas
situés dans les zones rurales ou de montagne programment davantage de
films Art et Essai. On peut y voir l'effet des politiques publiques de
soutien au cinéma culturel en dehors des grands centres urbains. Ce sont
souvent des salles classées Art et Essai, parfois associatives, qui ont
un rôle important dans l'animation culturelle locale. C'est aussi
peut-être une réponse à une demande différente : dans des zones moins
concurrentielles, ces salles peuvent se permettre une programmation plus
engagée ou originale.
\end{quote}

\section{Conclusion}\label{conclusion}

Cette étude descriptive permet de dresser un panorama des cinémas
français :

\begin{itemize}
\tightlist
\item
  une répartition inégale mais active sur le territoire,
\item
  une forte prédominance des petits établissements,
\item
  des multiplexes qui concentrent une large part de la fréquentation,
\item
  et une programmation qui valorise le cinéma national et le cinéma
  indépendant selon les zones.
\end{itemize}

Les questions posées ont trouvé des éléments de réponse grâce à des
analyses visuelles claires et comparatives. Pour aller plus loin, des
approches multivariées ou des comparaisons temporelles
(avant/après-Covid) pourraient enrichir cette étude.

\section{Annexe (optionnelle)}\label{annexe-optionnelle}

\begin{Shaded}
\begin{Highlighting}[]
\CommentTok{\# Exemple de fonction personnalisée}
\CommentTok{\# get\_top\_regions \textless{}{-} function(df, var, top = 5) \{}
\CommentTok{\#   df \%\textgreater{}\% group\_by(\textasciigrave{}région administrative\textasciigrave{}) \%\textgreater{}\%}
\CommentTok{\#     summarise(val = sum(\{\{ var \}\}, na.rm = TRUE)) \%\textgreater{}\%}
\CommentTok{\#     top\_n(top, val)}
\CommentTok{\# \}}
\end{Highlighting}
\end{Shaded}


\end{document}
